% Modernised reading — fonds Alexandre Grothendieck, shelfmark 47, pages 1-10.
%
% This is an INTERPRETATION, not a transcription. Everything the critical
% apparatus carried moves into footnotes. In any dispute of substance the
% transcription governs: whoever cites Grothendieck must cite batch-01.fr.tex.
%
% Pass: Fable 5.1 (claude-fable-5-1), 2026-09-03 — first pass, unchecked against
% the pages by a human. The folder was read whole; it holds one batch, pages
% 1-10, of which page 1 is a cover sheet in his hand and pages 2-10 one
% continuous run of mathematics. The skill admitted Fable 5.1 on 2026-09-03 and
% says the two permitted models are not measured against each other.
%
% Notation fixed for the whole document. Derived functors are written without
% his blackboard R and L: pr_{2*}, f^* denote the derived functors on
% D_parf, and (x) the derived tensor product, as is now usual; the one place
% where an underived operation is meant (tensoring with a line bundle) is the
% same either way. The Pontrjagin product is F * G = m_*(F ⊠ G), with m the
% group law and ⊠ the external tensor product over S. The Weil sheaf — the
% Poincaré bundle, in today's name — on A x B is L or W as on the page; the
% sheaf on B x A obtained by swapping the factors is written L^s (his
% superscript s before the L). The top exterior power he writes with a hat is
% written Λ^n; t_A is the Lie algebra of A (the tangent sheaf along the zero
% section, a vector bundle on S), τ_A = Λ^n t_A its determinant, and
% ω_A = Λ^n t_A^∨ = τ_A^{-1}. His "Cw" with a raised dot and subscript Q, read
% as glyphs by the transcription, is the Chow ring: CH^·(A)_Q throughout.
%
% Substantive departures from the pages, each footnoted where it occurs:
% — page 8: the expansion of the total Chern class c(φ_L(cl L)) = (1 - D'/χ)^χ
%   carries χ^{k-1} in the denominator of the k-th term, and the divisibility
%   line under it χ^{k-2}; the Chern character line two lines above, which is
%   what the formula is expanded from, forces χ^k. The reading takes χ^k.
% — page 8: the sign in the relation between the polarisation δ of L and the
%   polarisation δ' of det Φ_L(L) is read with doubt on the page and depends
%   on how A is identified with the dual of B; the reading gives the relation
%   in Mukai's normalisation, where the composite is +χ(δ).
% — page 10, remark 4: the duality formula is written without the symmetry
%   s_A; as written it fails on a skyscraper e_s. The reading inserts s_A^*.
% — page 10, remark 3: the second Chow formula carries a single Todd class
%   T = Todd(t_A) where the derivation from the page's own (*) gives
%   Todd(t_A)·Todd(t_B); the two agree over a field, which is Kleiman's case.
% — page 4: "il suffit de définir un isom. Φ_{sL}Φ_L = id" is the page's
%   shorthand for the functor announced in the theorem, s_A^* ⊗ τ_A[-n], which
%   is what the computation that follows produces.
%
% What the folder announces and never establishes: Lemma 2 (Mumford) is
% admitted, not proved; the last step of Lemma 1 — the compatibility check at
% points of the zero section — ends in dots; the marginal generalisations of
% Proposition 1 to group schemes and to torsors are stated and not pursued; the
% question whether the Chern-class formulas of remark 2 hold modulo torsion is
% left open (with a question mark of his own), as is the claim that Kleiman's
% formulas miss a sign.
\documentclass[11pt,a4paper]{article}
% Le préambule commun à toutes les transcriptions du fonds Grothendieck.
%
% Il tient en une idée : **l'incertitude se compose, elle ne se résout pas.**
% Une transcription qui lisse un mot douteux a détruit la seule chose pour
% laquelle on la faisait — un lecteur doit pouvoir savoir, à chaque endroit,
% ce qui était sur le feuillet et ce qui a été ajouté. Les six macros qui
% suivent sont donc l'appareil critique, et non de la décoration.
%
% Les mêmes commandes sont reconnues par `scripts/render.mjs`, qui produit la
% vue de lecture du site. Une macro ajoutée ici doit l'être aussi là-bas,
% faute de quoi le rendu échouera bruyamment — ce qui est voulu : un rendu
% silencieusement faux est pire qu'un rendu absent.

\ProvidesPackage{grothendieck}[2026/08/08 Transcriptions du fonds Grothendieck]

\usepackage{amsmath,amssymb,amsthm}
\usepackage{mathrsfs}
\usepackage{stmaryrd}
% Les diagrammes commutatifs sont partout dans ce fonds : c'est la langue
% même de la géométrie algébrique de Grothendieck, et une transcription qui
% les rendrait en prose ne transcrirait rien.
\usepackage{tikz-cd}
% Français par défaut — c'est la langue des feuillets — mais l'anglais est
% chargé aussi : la traduction et le résumé sont des documents anglais, et la
% typographie française (espaces avant les deux-points) y serait une faute.
% Ces documents basculent par \selectlanguage{english} en tête de corps.
\usepackage[english,french]{babel}
\usepackage{fontspec}
\usepackage{geometry}
\geometry{a4paper,margin=25mm,marginparwidth=28mm}
\usepackage{marginnote}
\usepackage{xcolor}
% ulem, not soul. `soul`'s \st and \ul are built on a character-by-character
% scan that XeLaTeX's font machinery breaks: the document fails with an
% undefined control sequence rather than degrading. ulem does the same two jobs
% and is Unicode-engine safe. `normalem` stops it hijacking \emph, which the
% transcriptions use for Grothendieck's own emphasis.
\usepackage[normalem]{ulem}
\usepackage{hyperref}
\hypersetup{colorlinks=true,linkcolor=black,urlcolor=black,pdfborder={0 0 0}}

% --- Métadonnées du lot -----------------------------------------------------
% Reprises telles quelles par le script de rendu, qui les affiche en tête de
% la vue de lecture. Elles fixent ce que le fichier prétend couvrir : sans
% elles, une transcription flotte sans point d'attache dans un fonds de seize
% mille feuillets.

\newcommand{\folder}[1]{\gdef\grothFolder{#1}}
\newcommand{\batch}[1]{\gdef\grothBatch{#1}}
\newcommand{\leaves}[2]{\gdef\grothFirst{#1}\gdef\grothLast{#2}}
\newcommand{\foldertitle}[1]{\gdef\grothFoldertitle{#1}}
\gdef\grothFolder{?}\gdef\grothBatch{?}\gdef\grothFirst{?}\gdef\grothLast{?}\gdef\grothFoldertitle{}

% --- Le résumé --------------------------------------------------------------
%
% Il ouvre la lecture modernisée, sous le titre « Résumé », et il est composé
% en petit. Ce n'est pas un ornement : c'est le seul passage du document qui
% ne soit pas des mathématiques, et un lecteur doit pouvoir le reconnaître
% avant de l'avoir lu — puis le sauter s'il connaît déjà le sujet, ou s'y
% arrêter s'il ne le connaît pas. Le filet qui le suit marque la reprise du
% corps.

\newenvironment{resume}
  {\small\setlength{\parskip}{0.45em}}
  {\par\medskip\hrule\medskip}

% --- Mots-clés ---------------------------------------------------------------
%
% En anglais, en clôture du résumé de la lecture modernisée : le vocabulaire
% moderne sous lequel chercher ce que les feuillets construisent. Le manifeste
% du site (`npm run manifest`) les extrait pour étiqueter la cote entière —
% cette ligne est la source unique des tags, il n'y a pas de fichier à côté.
\newcommand{\keywords}[1]{\par\smallskip\noindent
  {\footnotesize\textsc{Keywords} --- \textit{#1}}\par}

% --- L'appareil critique ----------------------------------------------------

% Le numéro de feuillet, en marge. C'est l'ancre qui permet de régler un
% désaccord sur une formule en regardant *un* feuillet plutôt qu'une chemise
% de six cents. Le site s'en sert aussi pour tourner les pages du fac-similé
% quand on fait défiler la transcription.
\newcommand{\leaf}[1]{%
  \marginnote{\footnotesize\color{gray!70}\textsf{#1}}%
  \label{leaf:#1}\ignorespaces}

% Illisible : on ne devine pas. Un mot inventé qui se lit comme les autres est
% le pire résultat possible.
\newcommand{\ill}{\textcolor{red!60!black}{[\,\ldots\,]}}

% Lecture douteuse : le mot est donné, et signalé comme incertain.
\newcommand{\uncertain}[1]{\uline{#1}}

% Ajout de l'éditeur — une lettre manquante, un mot sous-entendu.
\newcommand{\add}[1]{\textcolor{blue!50!black}{[#1]}}

% Biffé par Grothendieck lui-même. Ce qu'il a rayé fait partie du document :
% une rature dit souvent où la pensée a changé de direction.
\newcommand{\struck}[1]{\sout{#1}}

% Note du transcripteur, sur l'état matériel du feuillet.
\newcommand{\note}[1]{\par\smallskip{\small\color{gray!60!black}\textit{[#1]}}\par\smallskip}

% Note marginale de Grothendieck, distincte du corps du texte.
\newcommand{\marginal}[1]{\par\smallskip\noindent\hspace{1em}%
  {\small\color{gray!60!black}#1}\par\smallskip}

% --- Titre ------------------------------------------------------------------

\AtBeginDocument{%
  \begin{center}
    {\large\bfseries Fonds Alexandre Grothendieck --- cote n\textsuperscript{o}~\grothFolder}\\[2pt]
    {\itshape\grothFoldertitle}\\[4pt]
    {\small feuillets \grothFirst--\grothLast\ (lot \grothBatch)}\\[2pt]
    {\footnotesize\color{gray!60!black}Transcription établie d'après le fac-similé numérisé
     par l'Université de Montpellier.\\
     Les lectures incertaines sont soulignées, les illisibles notées [\,\ldots\,],
     les ajouts entre crochets.}
  \end{center}
  \vspace{1em}}

\endinput


\folder{47}
\pages{1}{10}
\dating{s.d. — le groupe « Variétés abéliennes » (45 à 56) est daté 1961-1973}
\watermark{Édition de démonstration\\interprétation personnelle de l'œuvre}
\foldertitle{Produit de convolution sur les VA et transformation de Fourier dans les VA : notes manuscrites (s.d.) — lecture modernisée du dossier entier}

\begin{document}

\section*{Résumé}

\begin{resume}
Ces neuf feuillets construisent une transformation de Fourier pour les
variétés abéliennes, et démontrent qu'elle s'inverse.

La transformation de Fourier ordinaire prend une fonction sur la droite
réelle et la décompose en ondes : à chaque fréquence $\xi$ elle associe
l'intégrale de $f(x)e^{-ix\xi}$. Elle doit ses vertus à deux faits. Le
premier est qu'elle échange deux opérations de nature différente : le
\emph{produit de convolution} $f * g$, qui mélange deux fonctions en
déplaçant l'une le long de l'autre, devient le produit ordinaire
$\hat f \cdot \hat g$. Le second est qu'elle s'inverse : appliquée deux fois
elle rend la fonction de départ, retournée, $f(-x)$, à une constante près.

Une variété abélienne est un espace algébrique compact muni d'une addition —
la généralisation en toute dimension d'une courbe elliptique, dont on sait
additionner les points. Elle a une \emph{duale} $B$, qui paramètre ses
fibrés en droites de degré zéro : un point $b$ de $B$ est un fibré $L_b$ sur
$A$, et l'ensemble de ces fibrés forme un fibré en droites universel $L$ sur
le produit $A \times B$, le faisceau de Weil, aujourd'hui le fibré de
Poincaré. C'est lui qui tient le rôle de l'exponentielle $e^{-ix\xi}$ : un
« point » $b$ de la duale est une fréquence, et $L_b$ l'onde correspondante.
Les « fonctions » à transformer ne sont pas des fonctions mais des
\emph{faisceaux}, ou plus exactement des complexes de faisceaux ; et
l'intégrale est remplacée par la cohomologie le long des fibres. La
transformation est donc $\Phi_L(F) = pr_{2*}(pr_1^{*}F \otimes L)$ : on tire
$F$ sur $A \times B$, on le multiplie par l'onde universelle, on intègre le
long de $A$.

Les deux vertus subsistent, et ce sont les deux résultats du dossier. La
proposition 1 dit que $\Phi_L$ transforme le produit de convolution des
faisceaux — le \emph{produit de Pontrjagin}, obtenu en poussant $F \boxtimes G$
le long de l'addition $A \times A \to A$ — en produit tensoriel. Le théorème
dit que, si l'on transforme de $A$ vers $B$ puis de $B$ vers $A$ avec le
même noyau, on retrouve le faisceau de départ retourné par la symétrie
$a \mapsto -a$, décalé de $n$ degrés et tordu par un fibré en droites
canonique sur la base : le déterminant de l'algèbre de Lie. En particulier
$\Phi_L$ est une équivalence de catégories. La preuve ramène tout à un calcul,
la cohomologie du fibré de Poincaré le long de l'une des projections, qui se
concentre en un point : le lemme d'annulation de Mumford — un fibré de degré
zéro non trivial n'a aucune cohomologie — est admis, et le reste se déduit.

Les remarques qui suivent tirent les conséquences. Le rang et la
caractéristique d'Euler s'échangent ; les faisceaux ponctuels deviennent des
fibrés en droites, et les fibrés en droites de degré zéro des faisceaux
ponctuels. La transformée d'un fibré en droites de caractéristique $\chi$ se
comporte, heuristiquement, comme $\chi$ copies d'une racine $\chi$-ième d'un
fibré sur la duale, ce dont il tire une formule pour ses classes de Chern et
une conjecture de divisibilité. Enfin, en passant aux classes de cycles par le
théorème de Riemann--Roch, il obtient une transformation de Fourier sur
l'anneau de Chow rationnel, échangeant encore convolution et produit — au
signe $(-1)^n$ près, dont il note que les formules publiées par Kleiman l'ont
omis.

Tout est écrit sur une base quelconque, pour des \emph{schémas abéliens},
c'est-à-dire des familles de variétés abéliennes ; c'est ce qui fait
apparaître la torsion par le déterminant de l'algèbre de Lie et les classes de
Todd des remarques finales, qui disparaissent sur un corps. Ce qui reste
alors est, à la lettre, le théorème d'inversion de Mukai (1981) et la
transformation de Fourier de Beauville sur l'anneau de Chow (1983). Le
dossier, non daté, cite les \emph{Dix exposés sur la cohomologie des schémas}
de 1968 ; il leur est donc postérieur, et antérieur d'une dizaine d'années au
moins à ces deux articles, qu'il ne cite évidemment pas.

\keywords{Fourier-Mukai transform, abelian scheme, dual abelian scheme, Poincaré bundle, Pontrjagin product, derived category of perfect complexes, Mukai inversion theorem, Mumford vanishing theorem, theorem of the square, Künneth formula, composition of kernels, Grothendieck-Riemann-Roch, Todd class, Chern character, Fourier transform on Chow groups, Néron-Severi group, polarisation, determinant of cohomology, Grothendieck duality, contragredient automorphism}
\end{resume}

\pagerange{1}{10}
\section*{Le fil du dossier, et les conventions}

Le dossier tient en dix feuillets. Le premier est une chemise qui porte, de
sa main, le titre que l'inventaire reproduit : « Produit de convolution sur
les VA et transformation de Fourier dans les VA », où VA abrège variétés
abéliennes. Les neuf autres sont une suite continue à l'encre bleue,
paginée 2 à 10 par les archivistes, sans un feuillet blanc ni une page
sautée. Le fil est celui-ci :

\begin{itemize}
  \item \emph{pages 2 et 3} : la définition du foncteur $\Phi_L$ et la
  proposition 1, qui échange convolution et produit tensoriel ; la
  démonstration ramène au faisceau de Weil et n'emploie que le changement de
  base, la formule de projection, le théorème du carré et Künneth ;
  \item \emph{pages 3 à 5} : le théorème d'inversion — d'abord écrit
  « Proposition 2 », et cité ainsi dans la suite —, son corollaire, et sa
  démonstration, qui calcule le noyau du foncteur composé et le ramène,
  par le lemme 1 et son corollaire, à la cohomologie du faisceau de Weil ;
  \item \emph{pages 6 et 7} : le lemme 2, d'annulation, attribué à Mumford et
  admis, et la déduction du lemme 1 à partir de lui, dont la dernière
  vérification est laissée en pointillé ;
  \item \emph{pages 7 à 10} : cinq remarques — images des faisceaux
  élémentaires et échange du rang et de la caractéristique d'Euler ; l'image
  d'un faisceau inversible, ses classes de Chern et une question de
  divisibilité ; la version dans l'anneau de Chow, par Riemann--Roch ;
  la dualité ; le transport de structure.
\end{itemize}

\subsection*{Les conventions, fixées ici pour tout ce qui suit}

$S$ est un schéma quelconque, $A$ et $B$ des schémas abéliens sur $S$ de
dimension relative $n$, $g_A \colon A \to S$ le morphisme structural,
$e_A \colon S \to A$ la section nulle, $m_A \colon A \times_S A \to A$
l'addition (le manuscrit écrit $\pi$ ou $\pi_A$), $s_A \colon A \to A$ la
symétrie $a \mapsto -a$. Tous les produits sont pris sur $S$, et l'indice est
omis quand il n'y a pas d'ambiguïté.

$D_{\mathrm{parf}}(A)$ est la catégorie dérivée des complexes parfaits sur
$A$. Tous les foncteurs $f^{*}$, $f_{*}$, $\otimes$ sont \emph{dérivés} ; le
manuscrit les affuble d'un $\mathbb{R}$ ou d'un $\mathbb{L}$ et d'un $L$ sur
le signe $\otimes$ et sur le signe $*$, que la lecture omet, l'usage ayant
changé.\footnote{Une note marginale de la page 4 dit qu'il est « inutile de
se borner aux complexes parfaits » et propose $D^{b}_{\mathrm{qcoh}}$ ; la
lecture reste dans $D_{\mathrm{parf}}$, où tous les foncteurs employés
(images directes par des morphismes propres et lisses, images réciproques,
produits tensoriels) préservent la perfection sans hypothèse
supplémentaire.} Pour $F$ sur $A$ et $G$ sur $A'$, le produit tensoriel
externe est $F \boxtimes G = pr_1^{*}F \otimes pr_2^{*}G$ sur $A \times A'$ ;
le manuscrit l'écrit $F \otimes_S G$.

Le \emph{produit de Pontrjagin} de $F, G \in D_{\mathrm{parf}}(A)$ est
\[
  F * G = m_{A*}(F \boxtimes G).
\]
Il est associatif et commutatif à isomorphisme canonique près, et le
faisceau $e_{A*}\mathcal{O}_S$ — écrit $e_A$ — en est l'unité ; pour une
section $s$ de $A$, le faisceau $e_s = s_{*}\mathcal{O}_S$ vérifie
$e_s * F = t_{s*}F$, translaté de $F$ par $s$.

Le \emph{schéma abélien dual} de $A$ est $A^{*}$ ; le \emph{faisceau de
Weil} sur $A \times A^{*}$ — aujourd'hui le \emph{fibré de Poincaré} — est le
module inversible universel, trivialisé le long de $e_A \times A^{*}$ et de
$A \times e_{A^{*}}$ (\emph{birigidifié}), dont la restriction à $A \times
\{b\}$ est le fibré algébriquement équivalent à zéro que le point $b$
représente. Un module inversible birigidifié $L$ sur $A \times B$
\emph{définit une dualité} entre $A$ et $B$ si l'homomorphisme
$\varphi \colon A \to B^{*}$ qu'il induit est un isomorphisme ; $L$ est alors
le faisceau de Weil transporté par $\varphi$, et le dossier écrit
indifféremment $L$ ou $W$. Le module sur $B \times A$ déduit de $L$ par
l'échange des facteurs est ${}^{s}L$.

Le \emph{noyau} $L$ étant fixé, la \emph{transformation de Fourier} du dossier
est
\[
  \Phi_L \colon D_{\mathrm{parf}}(A) \longrightarrow D_{\mathrm{parf}}(B),
  \qquad
  \Phi_L(F) = pr_{2*}\bigl(pr_1^{*}F \otimes L\bigr),
\]
et $\varphi_L \colon K(A) \to K(B)$ l'application qu'elle induit sur les
groupes de Grothendieck des complexes parfaits. C'est ce que la littérature
appelle depuis 1981 la \emph{transformation de Fourier--Mukai} de noyau
$L$.\footnote{S. Mukai, \emph{Duality between $D(X)$ and $D(\hat X)$ with its
application to Picard sheaves}, Nagoya Math.\ J.\ 81 (1981), sur un corps ;
la version sur une base quelconque, avec la torsion par le déterminant de
l'algèbre de Lie qu'on trouvera au théorème, est celle des traitements
relatifs ultérieurs, par exemple G.~Laumon, \emph{Transformation de Fourier
généralisée} (1996). Le dossier ne pouvait citer ni l'un ni l'autre ; ce
qu'il cite, ce sont les \emph{Dix exposés} de 1968 et « les résultats de
Mumford », c'est-à-dire les \emph{Abelian varieties} de 1970 ou le cours qui
les a précédées.}

Enfin $t_A = e_A^{*}T_{A/S}$ est l'\emph{algèbre de Lie} de $A$, un module
localement libre de rang $n$ sur $S$ — le manuscrit dit « faisceau tangent le
long de la section nulle » —, $\tau_A = \Lambda^{n}t_A$ son déterminant et
$\omega_A = \Lambda^{n}t_A^{\vee} = \tau_A^{-1}$.\footnote{Le manuscrit note
la puissance extérieure maximale $\hat\Lambda$, avec un accent, puis
$\hat\Lambda^{n}$ ; il pose $\tau_A \simeq \hat\Lambda\, t_A$ et
$\omega_A = \hat\Lambda\, t_A^{\vee}$. Ces modules sur $S$ sont tirés sur $A$
ou sur $B$ sans changement de notation.} L'anneau de Chow rationnel est
$\mathrm{CH}^{\cdot}(A)_{\mathbf{Q}}$.\footnote{Le manuscrit écrit un groupe
« Cw » — un $C$, une lettre basse, un point en exposant et l'indice
$\mathbf{Q}$ — que la transcription rend par les glyphes lus. La page 10
l'oppose à une « théorie cohomologique », et les formules qui l'emploient
sont celles de l'anneau de Chow à coefficients rationnels ; c'est ainsi que
la lecture l'écrit.}

\subsection*{Ce que le dossier annonce et n'établit pas}

Le lemme 2, d'annulation, est attribué à Mumford et admis. La démonstration
du lemme 1 s'arrête, à la page 7, sur une « vérification de compatibilité »
en pointillé. La note marginale de la page 2 esquisse deux généralisations de
la proposition 1 — aux schémas en groupes munis d'un module $L$
multiplicatif, et aux torseurs — sans les poursuivre. La question de la page
8, si les formules de classes de Chern valent modulo torsion, porte son point
d'interrogation. Enfin, l'affirmation que Kleiman « n'a pas tenu compte des
facteurs $(-1)^n$ » est la sienne, et la lecture la rapporte sans l'avoir
vérifiée sur le texte de Kleiman.

\pagerange{2}{3}
\section*{La transformation échange convolution et produit tensoriel (pages 2 et 3)}

\subsection*{L'énoncé}

La proposition 1 est énoncée dans une généralité plus grande que celle du
théorème : le noyau n'est pas encore le faisceau de Weil, mais une famille
quelconque de fibrés de degré zéro.

\textbf{Proposition 1.} Soient $S$ un schéma, $A$ un schéma abélien sur $S$,
$X$ un $S$-schéma, et $L$ un module inversible sur $A \times_S X$, trivial
le long de $e_A \times_S X$, et tel que pour tout point $x$ de $X$,
d'image $s$ dans $S$, le fibré $L_x$ sur la fibre $A_x = A_s \otimes_{k(s)}
k(x)$ soit algébriquement équivalent à zéro. Alors le foncteur
\[
  \Phi_L \colon D_{\mathrm{parf}}(A) \longrightarrow D_{\mathrm{parf}}(X),
  \qquad \Phi_L(F) = pr_{2*}\bigl(pr_1^{*}F \otimes L\bigr)
\]
transforme le produit de Pontrjagin en produit tensoriel : on a un
isomorphisme de foncteurs triangulés
\[
  \Phi_L(F * G) \;\simeq\; \Phi_L(F) \otimes \Phi_L(G)
  \qquad (F, G \in D_{\mathrm{parf}}(A)),
\]
et par suite $\varphi_L \colon K(A) \to K(X)$ est un homomorphisme d'anneaux
de $(K(A), *)$ vers $(K(X), \cdot)$.\footnote{La page énonce d'abord la
version en $K$-théorie — « $\varphi_L$ transforme produit de Pontrjagin en
produit ordinaire » — puis « plus précisément » l'isomorphisme de foncteurs.
La première rédaction, biffée, parlait d'une « application additive » avant
que le foncteur ne soit nommé. Une note marginale, écrite en biais et
partiellement illisible, indique deux généralisations : à un schéma en
groupes $A$ quelconque muni d'un $L$ \emph{multiplicatif}, c'est-à-dire tel
que $(m \times \mathrm{id}_X)^{*}L \simeq pr_{13}^{*}L \otimes pr_{23}^{*}L$ —
c'est la seule propriété de $L$ que la démonstration utilise —, et à des
complexes sur des torseurs $P$, $Q$, $R$ sous $A$ munis d'une loi
$P \times Q \to R$, avec trois noyaux $L$, $L'$, $L''$ liés par la relation
analogue. Ni l'une ni l'autre n'est développée.}

\subsection*{La réduction au faisceau de Weil}

Le noyau se transporte par changement de base : pour un $S$-morphisme
$f \colon X' \to X$ et $L' = (\mathrm{id}_A \times f)^{*}L$, on a
$\Phi_{L'} \simeq f^{*} \circ \Phi_L$, puisque $pr_2 \colon A \times X \to X$
est propre et plat et que l'image directe dérivée commute alors à tout
changement de base ; comme $f^{*}$ commute au produit tensoriel, l'énoncé
pour $(A, X, L)$ entraîne l'énoncé pour $(A, X', L')$.

Or l'hypothèse faite sur $L$ — trivial le long de la section nulle,
algébriquement équivalent à zéro fibre par fibre — dit exactement que $L$ est
l'image réciproque du faisceau de Weil $L_0$ sur $A \times B$, $B = A^{*}$,
par un unique $S$-morphisme $X \to B$ : c'est la propriété universelle du
schéma abélien dual.\footnote{La page dit « l'hypothèse sur $L$ implique que
$L$ est $\simeq$ l'image inverse de $L_0$ par un morphisme $X \to B$ ». La
trivialisation le long de $e_A \times X$ est ce qui fixe l'isomorphisme et non
seulement la classe ; c'est pourquoi elle figure dans l'hypothèse.} On est
donc ramené au cas $X = B$, $L = L_0$, et c'est dans ce cas que la
démonstration est écrite.

\subsection*{La démonstration}

Elle repose sur quatre outils, tous formels sauf le troisième : le
changement de base pour les images directes le long d'un carré cartésien
dont une flèche est plate, la formule de projection, le \emph{théorème du
carré}, et la formule de Künneth. Les deux diagrammes de la page 3 fixent les
noms :
\[
\begin{tikzcd}[column sep=small]
A \times A \arrow[d, "m"'] & A \times A \times B \arrow[l, "pr_{12}"'] \arrow[d, "m \times \mathrm{id}_B"] \arrow[dr, "pr_3"] & \\
A & A \times B \arrow[l, "pr_1"] \arrow[r, "pr_2"'] & B
\end{tikzcd}
\qquad
\begin{tikzcd}[column sep=small]
A \times A \times B \arrow[r, "pr_3"] \arrow[d, "f"'] & B \arrow[d, "\Delta_B"] \\
A \times A \times B \times B \arrow[r, "pr_{34}"] & B \times B
\end{tikzcd}
\]
où le carré de gauche du premier diagramme est cartésien, et
$f = \mathrm{id}_{A \times A} \times \Delta_B$ rend le second
cartésien.\footnote{Le second diagramme de la page porte aussi une flèche
$pr'_{12}$ de $A \times A \times B \times B$ vers $A \times A$ et une flèche
courbe $pr_{12}$ ; elles servent à nommer les projections de la formule
ci-dessous et sont omises du dessin.} Les projections
$pr'_{12}, pr'_{13}, pr'_{24}$ sont celles de $A \times A \times B \times B$,
et $pr_{13} = pr_1 \times \mathrm{id}_B$, $pr_{23} = pr_2 \times
\mathrm{id}_B$ celles de $A \times A \times B$ vers $A \times B$.

Le \emph{théorème du carré}, sous la forme où il sert ici, dit que le
faisceau de Weil est multiplicatif dans la variable $A$ :
\[
  (m \times \mathrm{id}_B)^{*}L \;\simeq\; pr_{13}^{*}L \otimes pr_{23}^{*}L
  \quad \text{sur } A \times A \times B.
\]
Restreint à $A \times A \times \{b\}$ c'est l'énoncé $m^{*}L_b \simeq
pr_1^{*}L_b \otimes pr_2^{*}L_b$ qui caractérise les fibrés algébriquement
équivalents à zéro ; l'isomorphisme est fixé, et global, par la
birigidification.\footnote{La page écrit « th.\ des carrés ». C'est le
théorème du carré de Weil, dans la forme que lui donne Mumford
(\emph{Abelian varieties}, § 6 et § 8) : $L$ est dans $\mathrm{Pic}^{0}$ si
et seulement si $m^{*}L \simeq L \boxtimes L$. Dans le langage d'aujourd'hui,
la formule dit que le fibré de Poincaré est une \emph{biextension} de
$(A, B)$ par $\mathbf{G}_m$ ; le mot est de SGA 7, postérieur à Weil et
peut-être à ces pages.}

Soient alors $F, G \in D_{\mathrm{parf}}(A)$. Par définition, puis par
changement de base le long du carré cartésien de gauche, formule de
projection et transitivité $pr_2 \circ (m \times \mathrm{id}_B) = pr_3$ :
\begin{align*}
  \Phi_L(F * G)
  &= pr_{2*}\bigl(pr_1^{*}\, m_{*}(F \boxtimes G) \otimes L\bigr) \\
  &\simeq pr_{2*}\bigl((m \times \mathrm{id}_B)_{*}\, pr_{12}^{*}(F \boxtimes G) \otimes L\bigr) \\
  &\simeq pr_{2*}(m \times \mathrm{id}_B)_{*}\bigl(pr_{12}^{*}(F \boxtimes G) \otimes (m \times \mathrm{id}_B)^{*}L\bigr) \\
  &\simeq pr_{3*}\bigl(pr_{12}^{*}(F \boxtimes G) \otimes (m \times \mathrm{id}_B)^{*}L\bigr).
\end{align*}
Le théorème du carré remplace le dernier facteur par
$pr_{13}^{*}L \otimes pr_{23}^{*}L$, et l'on réorganise les quatre facteurs
sur $A \times A \times B \times B$ : le complexe
$pr_{12}^{*}(F \boxtimes G) \otimes pr_{13}^{*}L \otimes pr_{23}^{*}L$ est
l'image réciproque par $f$ de
\[
  pr^{\prime *}_{12}(F \boxtimes G) \otimes pr^{\prime *}_{13}L \otimes pr^{\prime *}_{24}L
  \;\simeq\;
  pr^{\prime *}_{13}(pr_1^{*}F \otimes L) \otimes pr^{\prime *}_{24}(pr_1^{*}G \otimes L),
\]
la seconde égalité n'étant qu'une réassociation du produit tensoriel : le
premier facteur ne dépend que des coordonnées $1$ et $3$, le second que des
coordonnées $2$ et $4$. Le changement de base le long du second carré, puis
Künneth, donnent
\begin{align*}
  \Phi_L(F * G)
  &\simeq pr_{3*}\, f^{*}\bigl(pr^{\prime *}_{13}(pr_1^{*}F \otimes L) \otimes pr^{\prime *}_{24}(pr_1^{*}G \otimes L)\bigr) \\
  &\simeq \Delta_B^{*}\, pr_{34*}\bigl(pr^{\prime *}_{13}(pr_1^{*}F \otimes L) \otimes pr^{\prime *}_{24}(pr_1^{*}G \otimes L)\bigr) \\
  &\simeq \Delta_B^{*}\bigl(pr_{2*}(pr_1^{*}F \otimes L) \boxtimes pr_{2*}(pr_1^{*}G \otimes L)\bigr) \\
  &= \Delta_B^{*}\bigl(\Phi_L(F) \boxtimes \Phi_L(G)\bigr)
  \;\simeq\; \Phi_L(F) \otimes \Phi_L(G),
\end{align*}
ce qui est la proposition.\footnote{Sur la page, le « $\otimes L$ » de la
première ligne a d'abord été oublié et ajouté au-dessus, et l'avant-dernière
formule contient une parenthèse vide, $\mathbb{L}f^{*}(\quad)$, qu'un trait
relie à l'expression soulignée deux lignes plus haut ; la lecture l'a
remplie. La formule de Künneth utilisée est celle des complexes parfaits sur
un produit fibré de morphismes propres et plats, $pr_{34*}(P \boxtimes Q)
\simeq pr_{2*}P \boxtimes pr_{2*}Q$, qui résulte elle-même du changement de
base et de la formule de projection.}

\pagerange{3}{4}
\section*{Le théorème d'inversion (pages 3 et 4)}

\subsection*{L'énoncé}

Le noyau est désormais un module inversible définissant une dualité entre
$A$ et $B$ ; on transforme de $A$ vers $B$ avec $L$, puis de $B$ vers $A$
avec ${}^{s}L$.

\textbf{Théorème.} Soient $A$, $B$ deux schémas abéliens sur $S$ de
dimension relative $n$ et $L$ un module inversible birigidifié sur
$A \times B$ définissant une dualité entre $A$ et $B$. Le foncteur composé
\[
  D_{\mathrm{parf}}(A) \xrightarrow{\;\Phi_L\;} D_{\mathrm{parf}}(B)
  \xrightarrow{\;\Phi_{{}^{s}L}\;} D_{\mathrm{parf}}(A)
\]
est canoniquement isomorphe à
\[
  F \longmapsto s_A^{*}F \otimes \tau_A\,[-n],
\]
où $s_A^{*} \simeq s_{A*}$ est la symétrie et $\tau_A = \Lambda^{n}t_A$ le
déterminant de l'algèbre de Lie de $A$, tiré de $S$ ; et l'on a un
isomorphisme canonique $\tau_A \simeq \tau_B$. En particulier $\Phi_L$ et
$s_A^{*}\Phi_{{}^{s}L}$ sont des équivalences de catégories triangulées, et
sur les groupes de Grothendieck le composé
$K(A) \xrightarrow{\varphi_L} K(B) \xrightarrow{\varphi_{{}^{s}L}} K(A)$
est $(-1)^n\, \tau_A \cdot s_A^{*}$, de sorte que $\varphi_L$ et
$(-1)^n\,\omega_A \cdot s_A^{*}\varphi_{{}^{s}L}$ sont des isomorphismes
inverses l'un de l'autre.\footnote{L'énoncé s'appelait « Proposition 2 » et le
mot est biffé pour « Théorème » ; la suite du dossier continue de dire
« prop.\ 2 », et la lecture dit « le théorème ». La page énonce d'abord la
version en $K$-théorie, « le composé est $(-1)^n s_A^{*}$ », où la classe de
$\tau_A$ est sous-entendue — sur $S$ local, ou dès que $\tau_A$ est trivial,
elle vaut $1$ — puis, « plus précisément », le foncteur, avec la torsion. La
lecture met la torsion dans les deux. Sur un corps, $\tau_A$ est trivial et
l'énoncé est le théorème d'inversion de Mukai : $\Phi_{\hat P} \circ \Phi_P
\simeq (-1)^{*}[-g]$. L'isomorphisme $\tau_A \simeq \tau_B$ — le
déterminant de l'algèbre de Lie d'un schéma abélien est celui de son dual —
est affirmé sur la page sous la forme $\tau_A \simeq \hat\Lambda\, t_A
\simeq \hat\Lambda\, t_B$ ; il résulte du lemme 1 ci-dessous, appliqué aux
deux projections du faisceau de Weil, voir la note à ce lemme.}

\subsection*{Le corollaire : le produit tensoriel devient une convolution}

Appliquons la proposition 1 au noyau ${}^{s}L$ et aux complexes $\Phi_L(F)$,
$\Phi_L(G)$ sur $B$, puis le théorème deux fois :
\[
  \Phi_{{}^{s}L}\bigl(\Phi_L(F) * \Phi_L(G)\bigr)
  \simeq \Phi_{{}^{s}L}\Phi_L(F) \otimes \Phi_{{}^{s}L}\Phi_L(G)
  \simeq s_A^{*}(F \otimes G) \otimes \tau_A^{2}\,[-2n]
  \simeq \Phi_{{}^{s}L}\Phi_L(F \otimes G) \otimes \tau_A\,[-n].
\]
Comme $\Phi_{{}^{s}L}$ est une équivalence et commute au produit tensoriel
par un module tiré de $S$ et aux décalages, il vient :

\textbf{Corollaire.} Sous les hypothèses du théorème,
\[
  \Phi_L(F \otimes G) \;\simeq\; \bigl(\Phi_L(F) * \Phi_L(G)\bigr) \otimes \omega_A\,[n],
\]
ou, ce qui revient au même, le foncteur $F \mapsto \Phi_L(F \otimes
\omega_A[n]) \simeq \Phi_L(F) \otimes \omega_A[n]$ transforme le produit
tensoriel en produit de Pontrjagin. Sur les groupes de Grothendieck :
$\varphi_L(xy) = (-1)^n\,\omega_A\,\bigl(\varphi_L(x) * \varphi_L(y)\bigr)$.\footnote{Le début de l'énoncé est largement remanié sur la page ; on y lit
encore, biffés, un $\Phi_L^{\omega_A[n]}$, « $\overset{L}{\otimes}$ en
$\overset{L}{*}$ » et « $x \mapsto (-1)^n\,\mathrm{cl}\,\omega_A\,
\varphi_L(x)$ », et, non biffé, la formulation par $F \mapsto \Phi_L(F
\otimes \omega_A[n])$ que la lecture reprend. La page n'écrit pas la
déduction ; celle donnée ici est la seule que les deux énoncés précédents
permettent, et c'est aussi celle de Mukai (1981, 3.7). Dans la numérotation de
la page ce corollaire est « Corollaire 2 », le chiffre étant douteux.}

\pagerange{4}{5}
\section*{La démonstration du théorème : le noyau du composé (pages 4 et 5)}

\subsection*{Le composé de deux transformations à noyau}

Le composé de deux foncteurs à noyau est un foncteur à noyau, et son noyau
s'obtient en composant les noyaux. Ici
$\Phi_{{}^{s}L} \circ \Phi_L$ est de la forme
$F \mapsto pr_{2*}\bigl(pr_1^{*}F \otimes M\bigr)$ avec
\[
  M = pr_{13*}\bigl(pr_{12}^{*}L \otimes pr_{23}^{*}\,{}^{s}L\bigr)
  \;\in\; D_{\mathrm{parf}}(A \times A),
\]
les projections étant celles de $A \times B \times A$.\footnote{La page
écrit « $M \in D_{\mathrm{parf}}(A)$ » et nomme en marge l'espace
$A \times B \times A$ ; le noyau vit sur $A \times A$. Elle dit aussi qu'il
« suffit de définir un isomorphisme $\Phi_{{}^{s}L}\Phi_L =
\mathrm{id}_{D_{\mathrm{parf}}(A)}$ » ; c'est un raccourci pour le foncteur
annoncé, $s_A^{*}(-) \otimes \tau_A[-n]$, qui est ce que le calcul produit.}
Il s'agit donc de calculer $M$, et l'on va voir que c'est le faisceau
structural du graphe de la symétrie, décalé et tordu.

\subsection*{Le noyau comme image réciproque du faisceau de Weil}

Posons $C = A \times A$ et identifions $A \times B \times A$ à $C \times B$.
Le module inversible $N = pr_{12}^{*}L \otimes pr_{23}^{*}\,{}^{s}L$ sur
$C \times B$ a pour fibre en $((a, a'), b)$ le produit $L_{(a,b)} \otimes
L_{(a',b)}$ ; il est birigidifié de façon canonique et correspond donc à un
homomorphisme $\Psi \colon C \to B^{*}$, à savoir
\[
  \Psi = \varphi \circ m_A \colon A \times A \longrightarrow A \longrightarrow B^{*},
\]
où $\varphi \colon A \xrightarrow{\;\sim\;} B^{*}$ est l'isomorphisme défini
par $L$ : $\Psi$ est l'homomorphisme dont les deux composantes sont
$\varphi$.\footnote{« deux » est ajouté en interligne sur la page. Que $N$
corresponde à $\varphi \circ m_A$ est le théorème du carré de la section
précédente lu dans l'autre variable : $L_{(a,b)} \otimes L_{(a',b)} \simeq
L_{(a+a',\,b)}$.} Le carré
\[
\begin{tikzcd}
C \times B \arrow[r, "\Psi \times \mathrm{id}_B"] \arrow[d, "pr_1"'] & B^{*} \times B \arrow[d, "pr'_1"] \\
C \arrow[r, "\Psi"] & B^{*}
\end{tikzcd}
\]
est cartésien, et $N \simeq (\Psi \times \mathrm{id}_B)^{*}W_B$, où $W_B$
est le faisceau de Weil sur $B^{*} \times B$ ; par changement de base,
\[
  M = pr_{1*}(N) \;\simeq\; \Psi^{*}\bigl(pr'_{1*}W_B\bigr).
\]
La projection est ici celle sur le premier facteur.\footnote{La ligne de la page écrit $pr_2$, $pr'_2$ là où le diagramme, le
lemme et le corollaire qui suivent ont $pr_1$, $pr'_1$ : c'est la projection
sur le premier facteur qui est en cause, et la lecture la nomme ainsi
partout.} Tout est ainsi ramené à un calcul universel, la cohomologie du
faisceau de Weil le long d'une projection, qui fait l'objet du lemme 1.

\subsection*{Le lemme 1 et son corollaire}

\textbf{Lemme 1.} Soient $B$ un schéma abélien sur $S$ de dimension relative
$n$, $W$ le faisceau de Weil sur $B^{*} \times_S B$, $p \colon B^{*} \times_S
B \to B^{*}$ la première projection et $e \colon S \to B^{*}$ la section
nulle. On a un isomorphisme canonique
\[
  p_{*}W \;\simeq\; e_{*}\bigl(\tau_{B^{*}}\bigr)[-n],
  \qquad \tau_{B^{*}} = \Lambda^{n}t_{B^{*}},
\]
c'est-à-dire $p_{*}W \simeq e_{*}(\omega_{B^{*}}^{-1})[-n]$.\footnote{La
seconde forme est la note marginale de la page : « $\simeq e_{*}(\omega^{-1})$
où $\omega = \det(\Omega^{1}_{B^{*}/S})$ », le $\det$ étant douteux. Sur un
corps c'est le lemme central de Mukai (1981, § 2) : $p_{*}P \simeq
k(\hat 0)[-g]$. Le même calcul fait le long de la seconde projection donne
$e_{B*}(\tau_B)[-n]$ sur $B$, et les cohomologies totales
$R\Gamma(B^{*} \times B, W)$ calculées des deux façons sont $\tau_{B^{*}}[-n]$
et $\tau_B[-n]$ : d'où $\tau_{B^{*}} \simeq \tau_B$, l'isomorphisme que le
théorème affirme. La démonstration du lemme est aux pages 6 et 7.}

\textbf{Corollaire.} Soient $C$, $B$ des schémas abéliens sur $S$, $N$ un
module inversible birigidifié sur $C \times_S B$ et $\Psi \colon C \to B^{*}$
l'homomorphisme correspondant. Alors
\[
  pr_{1*}N \;\simeq\; \Psi^{*}\bigl(e_{*}\tau_{B^{*}}\bigr)[-n].
\]
En particulier, si $\Psi$ est un épimorphisme, donc plat, et si $K =
\operatorname{Ker}\Psi$, le sous-schéma en groupes $K$ de $C$ est plat et de
présentation finie sur $S$, son faisceau structural est un module parfait sur
$C$, et
\[
  pr_{1*}N \;\simeq\; \mathcal{O}_K \otimes \tau_{B^{*}}\,[-n].
\]
C'est le lemme 1 tiré en arrière par $\Psi$.\footnote{Les symboles entre la parenthèse et le $[-n]$ de la première formule
sont surchargés sur la page et lus avec doute ; le lemme 1 ne laisse qu'une
lecture possible. Pour la seconde : $\Psi$ étant plat, $\Psi^{*}e_{*}
\mathcal{O}_S$ est le faisceau structural de la fibre de $\Psi$ au-dessus de
la section nulle, qui est $K$, et $\Psi^{*}$ d'un module tiré de $S$ est ce
module. Que $\mathcal{O}_K$ soit parfait sur $C$ tient à ce que $K$ est plat
sur $S$ et $C$ lisse sur $S$ ; la page hésite — « défini », puis
« $\mathcal{O}_K$ », biffés — avant d'écrire « parfait », lu avec doute.}

\subsection*{Fin de la démonstration}

Dans le cas du théorème, $C = A \times A$ et $\Psi = \varphi \circ m_A$ avec
$\varphi$ un isomorphisme ; $\Psi$ est donc un épimorphisme plat, et son noyau
$K$ est le noyau de l'addition $m_A \colon A \times A \to A$, l'antidiagonale
$\{(a, -a)\}$, c'est-à-dire le \emph{graphe de la symétrie} $s_A$. Le
corollaire donne
\[
  M \;\simeq\; \mathcal{O}_{\Gamma_{s_A}} \otimes \tau_{B^{*}}\,[-n],
\]
et un foncteur à noyau $\mathcal{O}_{\Gamma_u}$, graphe d'un automorphisme
$u$, est $u_{*}$ ; ici $s_{A*} = s_A^{*}$. Enfin $\tau_{B^{*}} \simeq \tau_A$
par l'isomorphisme $\varphi$. Le théorème est donc démontré modulo le
lemme 1.\footnote{La page écrit « $\varphi \colon A \simeq C^{*}$ », la
lettre étant nette ; l'argument demande $B^{*}$, et c'est ce que la lecture
écrit. La conclusion est marquée « modulo dém.\ du Lemme 1 », lu avec doute.}

\pagerange{6}{7}
\section*{Le lemme d'annulation de Mumford et la cohomologie du faisceau de Weil (pages 6 et 7)}

\subsection*{Le lemme admis}

\textbf{Lemme 2 (Mumford).} Soient $B$ une variété abélienne sur un corps
algébriquement clos $k$ et $L$ un module inversible sur $B$ algébriquement
équivalent à zéro et non isomorphe à $\mathcal{O}_B$. Alors $H^{i}(B, L) = 0$
pour tout $i$.\footnote{La page écrit $H^{i}(X, L)$, la lettre $X$ pour $B$.
Un encadré entièrement biffé et le début d'une démonstration, biffé aussi —
« l'assertion sur les $H^{0}$ est triviale, et il reste à prouver que
$H^{i}(X, L) = 0$ » — précèdent la décision d'admettre le lemme. C'est le
théorème d'annulation de Mumford, \emph{Abelian varieties} (1970), § 8, dont
la démonstration est en quelques lignes : si $H^{i}(L) \neq 0$ pour un $i$
minimal, la factorisation $m^{*}L \simeq L \boxtimes L$ et Künneth donnent
$H^{i}(L) \simeq H^{0}(L) \otimes H^{i}(L)$, donc $H^{0}(L) \neq 0$, ce qui
force $L \simeq \mathcal{O}_B$ pour un fibré de degré zéro.}

« Nous allons ici admettre le lemme 2 », écrit-il, et montrer comment on en
déduit le lemme 1. L'énoncé est sur un corps ; il s'applique fibre à fibre,
et c'est ce qui permet de localiser.

\subsection*{La déduction du lemme 1}

On revient à $p \colon B^{*} \times_S B \to B^{*}$ et au faisceau de Weil $W$.
On écrit $e$ pour $e_{B^{*}}$, $g \colon B \to S$ pour le morphisme
structural et $e' = e \times \mathrm{id}_B \colon B \to B^{*} \times B$ ;
le carré
\[
\begin{tikzcd}
B^{*} \times_S B \arrow[d, "p"'] & B \arrow[l, "e'"'] \arrow[d, "g"] \\
B^{*} & S \arrow[l, "e"]
\end{tikzcd}
\]
est cartésien, et $e^{\prime *}W \simeq \mathcal{O}_B$ par la rigidification.

\emph{Le support.} Pour un point $b$ de $B^{*}$ hors de la section nulle, la
restriction de $W$ à la fibre $B_b$ est un fibré algébriquement équivalent à
zéro et non trivial ; par le lemme 2, appliqué après extension à une clôture
algébrique, toute sa cohomologie est nulle, et comme $p$ est propre et plat
la restriction dérivée de $p_{*}W$ en $b$ est nulle. Le complexe parfait
$p_{*}W$ est donc concentré sur la section nulle $e(S)$.

\emph{L'homomorphisme.} Comme $p$ a des fibres de dimension $n$, les
$R^{i}p_{*}W$ sont nuls pour $i > n$, et $p_{*}W$ s'envoie canoniquement sur
son dernier faisceau de cohomologie, $R^{n}p_{*}W\,[-n]$. Définir un
homomorphisme
\[
  p_{*}W \longrightarrow e_{*}(\mathcal{O}_S) \otimes_{\mathcal{O}_S} \tau_{B^{*}}\,[-n]
\]
revient donc à définir $R^{n}p_{*}W \to e_{*}\tau_{B^{*}}$, c'est-à-dire, par
adjonction, $e^{*}R^{n}p_{*}W \to \tau_{B^{*}}$. Or, par changement de base
le long du carré ci-dessus et parce que $e^{\prime *}W \simeq \mathcal{O}_B$,
\[
  e^{*}R^{n}p_{*}W \;\simeq\; R^{n}g_{*}(\mathcal{O}_B),
\]
et l'on sait que
\[
  R^{n}g_{*}(\mathcal{O}_B) \;\simeq\; \Lambda^{n}R^{1}g_{*}(\mathcal{O}_B)
  \;\simeq\; \Lambda^{n}t_{B^{*}} = \tau_{B^{*}},
\]
puisque $R^{1}g_{*}\mathcal{O}_B$ est l'algèbre de Lie du schéma abélien dual
$B^{*}$ et que la cohomologie de $\mathcal{O}_B$ est l'algèbre extérieure sur
son $H^{1}$.\footnote{Ce sont deux résultats de la théorie du schéma de
Picard : $\mathrm{Lie}(B^{*}) = R^{1}g_{*}\mathcal{O}_B$, et $R^{\cdot}
g_{*}\mathcal{O}_B \simeq \Lambda^{\cdot}R^{1}g_{*}\mathcal{O}_B$, qui pour un
schéma abélien vaut sur toute base (Mumford, § 13, ou Berthelot--Breen--Messing
pour la version relative). La page dit « on sait que l'on a un isom
canonique » et ne cite rien. Les deux $\otimes_S$ qui précèdent $\tau_{B^{*}}$
dans les formules de la page sont biffés.} On prend pour
$e^{*}R^{n}p_{*}W \to \tau_{B^{*}}$ l'isomorphisme identique ainsi obtenu.

\emph{L'isomorphisme.} « Je dis que l'homomorphisme précédent est un
isomorphisme. » Il suffit de le vérifier après restriction dérivée $i_b^{*}$
en tout point $b$ de $B^{*}$ — un morphisme de complexes parfaits qui est un
isomorphisme en chaque point l'est.\footnote{La page invoque une « histoire
de \ldots-cylindres », le premier mot illisible : le cône du morphisme est un
complexe parfait dont toutes les restrictions ponctuelles sont nulles, donc
nul, par le lemme de Nakayama appliqué à son dernier faisceau de cohomologie
non nul.} Aux points hors de $e(S)$, les deux membres sont nuls par le
lemme 2. En un point de $e(S)$, on se ramène par changement de base au cas où
$S$ est le spectre d'un corps algébriquement clos, et il reste « une
vérification de compatibilité », que la page laisse en pointillé.\footnote{Le
manuscrit s'arrête ici : « \ldots\ldots ». Ce qu'il reste à voir est que le
morphisme $p_{*}W \to k(0)[-n]$, sur une variété abélienne, est un
isomorphisme — autrement dit que $R^{n}p_{*}W$ est de longueur $1$ en $0$ et
non seulement supporté en $0$. C'est le contenu du lemme de Mukai (1981,
§ 2), qui l'obtient par dualité de Grothendieck et par le calcul de la
caractéristique d'Euler de $W$ sur $B^{*} \times B$ ; la lecture ne comble pas
le pointillé et se contente de dire que l'énoncé est vrai.}

\pagerange{7}{8}
\section*{Remarques 1 et 2 : les images élémentaires, et la transformée d'un faisceau inversible (pages 7 et 8)}

\subsection*{Remarque 1 : faisceaux ponctuels, fibrés de degré zéro, rang et caractéristique d'Euler}

Les hypothèses sont celles du théorème. Pour une section $s$ de $A$ on écrit
$e_s = s_{*}\mathcal{O}_S$, et $L_s = L|_{s \times B}$ le fibré
algébriquement équivalent à zéro sur $B$ que $s$ représente ; de même
$L_t = L|_{A \times t}$ pour une section $t$ de $B$.

\emph{a)} $\Phi_L(e_A) \simeq \mathcal{O}_B$ et $\Phi_L(e_s) \simeq L_s$ :
c'est la définition, puisque $pr_1^{*}e_s \otimes L$ est $L$ restreint à
$s \times B$. Comme $e_s * F$ est le translaté $t_{s*}F$, la proposition 1
donne
\[
  \Phi_L(t_{s*}F) \;\simeq\; L_s \otimes \Phi_L(F) :
\]
la translation devient la multiplication par un fibré de degré zéro.

\emph{b)} $\Phi_L(L_t) \simeq e_{-t} \otimes \tau_A\,[-n]$, ou encore
$\Phi_L(L_t \otimes \omega_A[n]) \simeq e_{-t}$ : la multiplication par un
fibré de degré zéro devient une translation, et le fibré trivial va sur le
faisceau ponctuel de l'origine. En effet, par le théorème du carré dans la
variable $B$, $pr_1^{*}L_t \otimes L \simeq (\mathrm{id}_A \times t_t)^{*}L$
où $t_t$ est la translation par $t$, de sorte que $\Phi_L(L_t) \simeq
t_t^{*}\,pr_{2*}L$, et $pr_{2*}L \simeq e_{B*}(\tau_B)[-n]$ par le lemme 1 ;
le translaté par $t$ du faisceau ponctuel de l'origine est celui de $-t$. De
même,
\[
  \Phi_L(F \otimes L_t) \;\simeq\; \Phi_L(F) * e_{-t}.
\]
La page hésite sur l'énoncé.\footnote{La page écrit d'abord « $\omega_A[n]$ en \ldots », biffé, avant de
fixer l'énoncé ; la lettre qui désigne le schéma dont $t$ est une section est
formée comme un $D$ et lue $B$. Ce sont les formules (3.1) de Mukai.}

\emph{c)} Notons $\varepsilon_A(F)$ le rang de la restriction dérivée
$e_A^{*}F$ à l'origine et $\chi_A(F)$ le rang de $g_{A*}F$, tous deux dans
$H^{0}(S, \mathbf{Z})$. Alors
\[
  \varepsilon_B(\Phi_L(F)) = \chi_A(F),
  \qquad
  \chi_B(\Phi_L(F)) = (-1)^n\,\varepsilon_A(F) :
\]
le rang à l'origine de la transformée est la caractéristique d'Euler, et la
caractéristique d'Euler de la transformée est le rang à l'origine, au signe
$(-1)^n$ près. La première égalité est le changement de base $e_B^{*}pr_{2*}
\simeq g_{A*}(\mathrm{id} \times e_B)^{*}$ et la rigidification ; la seconde
est $g_{B*}\Phi_L(F) \simeq g_{A*}(F \otimes pr_{1*}L) \simeq g_{A*}(F
\otimes e_{A*}\tau_A)[-n] \simeq e_A^{*}F \otimes \tau_A[-n]$, par le
lemme 1.\footnote{La définition de $\varepsilon$ est en marge et illisible,
sauf la mention « égalité dans $H^{0}(S, \mathbf{Z})$ » ; celle donnée ici
est la seule qui rende les deux égalités vraies, et elle est la nôtre.}

\subsection*{Remarque 2 : le déterminant de la transformée, et l'heuristique de la racine $\chi$-ième}

Le composé $\det \circ \Phi_L$ — le déterminant du complexe parfait
$\Phi_L(M)$ pour $M$ inversible sur $A$ — commute au changement de base et
définit donc un homomorphisme de foncteurs de Picard
\[
  \varphi \colon \underline{\mathrm{Pic}}_{A/S} \longrightarrow \underline{\mathrm{Pic}}_{B/S}.
\]
Soit $\lambda$ une section de $\underline{\mathrm{Pic}}_{A/S}$, de classe
$\delta$ dans le groupe de Néron--Severi $\underline{NS}_{A/S}$, définissant
un faisceau inversible que l'on notera encore $L$ sur $A$ ; soit $\chi =
\chi(\delta) \in H^{0}(S, \mathbf{Z})$ sa caractéristique d'Euler, supposée
partout non nulle. Alors $\varphi(\lambda)$ définit un faisceau inversible
$L'^{-1}$ sur $B$, de classe $-\delta'$ dans $NS(B/S)$, où $\delta'$ est
déterminée par $\delta$ par la relation
\[
  \tilde\delta' \circ \tilde\delta = \chi(\delta)\,\mathrm{id}_A,
\]
$\tilde\delta \colon A \to B$ et $\tilde\delta' \colon B \to A$ désignant les
homomorphismes associés aux polarisations.\footnote{La page écrit
$\tilde\delta'\tilde\delta = -\chi(\delta)\,\mathrm{id}_A$, et à côté une
ligne lue avec doute, « compatible avec $-(-\tilde\delta') = \tilde\delta'$ »,
qui suggère qu'il hésite sur le signe. Le signe dépend de la façon dont $A$
est identifié au dual de $B$ par $L$. Dans la normalisation de Mukai, où
$\varphi_L(a) = t_a^{*}L \otimes L^{-1}$ et où le fibré de Poincaré $P$ a
pour restriction à $A \times \{\alpha\}$ le fibré $P_\alpha$, sa proposition
3.11 donne $\varphi_L^{*}\det\Phi(L) \simeq L^{-\chi}$, d'où, $\varphi_L$
étant symétrique, $\varphi_{L'} \circ \varphi_L = \chi\,\mathrm{id}_A$ pour
$L' = \det\Phi(L)^{-1}$ : c'est la relation écrite ici, avec le signe $+$. La
lecture ne prétend pas que la page ait tort, seulement qu'elle a un autre
signe dans une autre convention, non explicitée.}

On peut se proposer de « calculer » la classe $\varphi_L(\mathrm{cl}\,L)$ dans
$K(B)$ en fonction des invariants de $\varphi(\lambda)$. Heuristiquement,
\[
  \varphi_L(\mathrm{cl}\,L) \;=\; \chi \cdot \text{« } L'^{-1/\chi} \text{ »} :
\]
la transformée d'un faisceau inversible de caractéristique $\chi$ se comporte
comme $\chi$ copies d'une racine $\chi$-ième de $L'^{-1}$. L'exposant est
entre guillemets sur la page, et l'heuristique est devenue un théorème :
$\Phi_L(L)$ est, à un décalage près, un fibré vectoriel de rang $|\chi|$ dont
l'image réciproque par $\tilde\delta$ est $\chi$ copies de $L^{-1}$, et
$\tilde\delta^{*}L' \simeq L^{\chi}$.\footnote{Mukai (1981), 3.11 : pour $L$
non dégénéré d'indice $i$, $\Phi(L)$ est concentré en degré $i$, et
$\varphi_L^{*}\Phi(L) \simeq H^{i}(A, L) \otimes L^{-1}$. Une note marginale
de la page dit « indépendant des résultats de Mumford » ; on comprend que
l'heuristique ne dépend pas du lemme 2 admis.}

Au niveau du caractère de Chern, l'heuristique s'écrit
\[
  \mathrm{ch}\,\varphi_L(\mathrm{cl}\,L)
  = \chi \exp\Bigl(-\frac{D'}{\chi}\Bigr)
  = \chi - D' + \frac{1}{2\chi}D'^{2} - \frac{1}{3!\,\chi^{2}}D'^{3} + \cdots
  + (-1)^n\frac{1}{n!\,\chi^{n-1}}D'^{n},
\]
où $D' \in \mathrm{CH}^{1}(B)_{\mathbf{Q}}$ est la classe de $L'$, et au
niveau de la classe de Chern totale
\[
  c\bigl(\varphi_L(\mathrm{cl}\,L)\bigr)
  = \Bigl(1 - \frac{D'}{\chi}\Bigr)^{\chi}
  = \sum_{k=0}^{n} (-1)^{k}\binom{\chi}{k}\frac{D'^{k}}{\chi^{k}}
  = 1 - D' + \frac{\chi(\chi-1)}{2\,\chi^{2}}D'^{2} - \cdots,
\]
la puissance $\chi$-ième étant développée par la formule du
binôme.\footnote{La page développe $(1 - D'/\chi)^{\chi}$ avec
$\chi^{k-1}$ au dénominateur du terme de degré $k$ — « $\frac{\chi(\chi-1)}
{2\chi}D'^{2}$ », « $\frac{\chi(\chi-1)(\chi-2)}{3!\,\chi^{2}}D'^{3}$ » —,
ce qui n'est pas le développement de ce qu'elle écrit à gauche, et ce qui
contredit la ligne du caractère de Chern : avec $c_1 = -D'$ et $c_2 =
\frac{\chi-1}{2}D'^{2}$ on aurait $\mathrm{ch}_2 = \frac{2-\chi}{2}D'^{2}$
au lieu de $\frac{1}{2\chi}D'^{2}$. La lecture prend $\chi^{k}$, qui est ce que
la ligne du $\mathrm{ch}$ et l'heuristique de la racine $\chi$-ième imposent
toutes deux. L'argument de $\varphi_L$ est surchargé sur la page dans les deux
formules.}

« Ces formules sont-elles bien vraies mod torsion ? » demande la page, avec
un point d'interrogation en marge. Elles sont vraies au moins sur un corps
algébriquement clos modulo équivalence homologique, où elles se lisent dans
la cohomologie. Elles donneraient alors des propriétés de divisibilité de
$\delta'$ : pour tout $k$,
\[
  \binom{\chi}{k}\frac{\delta'^{k}}{\chi^{k}}
  = \frac{\chi(\chi-1)\cdots(\chi-k+1)}{k!\,\chi^{k}}\,\delta'^{k}
\]
est « entier », c'est-à-dire provient d'un cycle à coefficients
entiers.\footnote{La page écrit $\chi^{k-2}$ au dénominateur, exposant qui
hérite du glissement de la formule précédente et s'en écarte encore d'un
cran ; la lecture écrit ce que l'intégralité des classes de Chern de
$\Phi_L(L)$ donne. La question « mod torsion » a aujourd'hui une réponse
positive sur un corps algébriquement clos : $\Phi_L(L)$ est un fibré, sa
classe de Chern totale est entière, son image réciproque par $\tilde\delta$
est $(1 - D)^{\chi}$ par le théorème de Mukai cité plus haut, et l'image
réciproque par une isogénie est injective sur l'anneau de Chow rationnel ;
l'égalité $c(\Phi_L(L)) = (1 - D'/\chi)^{\chi}$ vaut donc dans
$\mathrm{CH}^{\cdot}(B)_{\mathbf{Q}}$, et la divisibilité avec elle. Une
parenthèse de la page, « donc le $\neq$ car.\ $h$ », est lue avec doute et
n'est pas interprétée.}

\pagerange{9}{10}
\section*{Remarque 3 : la transformation de Fourier sur l'anneau de Chow (pages 9 et 10)}

On passe des complexes aux cycles par le théorème de Riemann--Roch de
Grothendieck, appliqué aux morphismes $m_A \colon A \times A \to A$ et
$pr_2 \colon A \times B \to B$. Ces deux morphismes sont lisses, et leur
fibré tangent relatif est tiré de $S$ : c'est $t_A$, l'algèbre de Lie. Leur
classe de Todd relative est donc la même classe
\[
  T = \mathrm{Todd}(t_A) \in \mathrm{CH}^{\cdot}(S)_{\mathbf{Q}},
\]
tirée sur $A$ ou sur $B$, et elle vaut $1$ dès que $t_A$ est trivial, par
exemple si $S$ est local — ou un corps.\footnote{La page dit « $\mathrm{Todd}
\,t_A = 1$, p.\ ex.\ si $t_A$ est trivial, p.\ ex.\ $S$ local ». Sur $A \times
A$ le fibré tangent est $pr_1^{*}T_A \oplus pr_2^{*}T_A$, chaque terme tiré de
$t_A$ ; le noyau de $dm_A$ est isomorphe à $t_A$ tiré de $S$.}

\emph{Le produit de Pontrjagin et le caractère de Chern.} Pour $x, y \in
K(A)$, Riemann--Roch pour $m_A$ donne
\[
  (*) \qquad \mathrm{ch}(x * y) = \bigl(\mathrm{ch}(x) * \mathrm{ch}(y)\bigr)\,T,
\]
où le produit de Pontrjagin des classes de cycles est $a * b = m_{A*}(a
\boxtimes b)$ ; en particulier $\mathrm{ch}$ est multiplicatif pour $*$ si
$T = 1$.

\emph{La transformation de Fourier des cycles.} Soit $D \in \mathrm{CH}^{1}(A
\times B)$ la classe de $L$, et
\[
  \Psi_D = \varphi_{\exp D} \colon
  \mathrm{CH}^{\cdot}(A)_{\mathbf{Q}} \longrightarrow \mathrm{CH}^{\cdot}(B)_{\mathbf{Q}},
  \qquad
  \Psi_D(a) = pr_{2*}\bigl(pr_1^{*}a \cdot \exp D\bigr).
\]
Riemann--Roch pour $pr_2$ donne alors, pour tout $x \in K(A)$,
\[
  \mathrm{ch}\,\varphi_L(x) = \Psi_D(\mathrm{ch}\,x)\,T .
\]
C'est la transformation de Fourier de l'anneau de Chow d'une variété
abélienne, telle que Beauville l'a définie et étudiée en 1983 à la suite de
Mukai ; la version cohomologique est de Lieberman (1968), et c'est elle que
Kleiman expose dans les \emph{Dix exposés}.\footnote{A.~Beauville,
\emph{Quelques remarques sur la transformation de Fourier dans l'anneau de
Chow d'une variété abélienne}, 1983 ; D.~Lieberman, \emph{Numerical and
homological equivalence of algebraic cycles on Hodge manifolds}, 1968 ;
S.~Kleiman, \emph{Algebraic cycles and the Weil conjectures}, dans \emph{Dix
exposés sur la cohomologie des schémas}, 1968. La page écrit « 10 Exposés
\ldots ». Le dossier, qui cite ce volume, lui est postérieur ; le nom
$\Psi_D$ et la définition par $\exp D$ sont les siens.}

\emph{Les deux formules.} Posons $a = \mathrm{ch}\,x$, $b = \mathrm{ch}\,y$.
En appliquant $\mathrm{ch}$ à la proposition 1, $\varphi_L(x * y) =
\varphi_L(x)\,\varphi_L(y)$, on trouve $\Psi_D(\mathrm{ch}(x * y))\,T =
\Psi_D(a)\,T\cdot\Psi_D(b)\,T$, puis, par $(*)$ et parce que $\Psi_D$ est
linéaire sur $\mathrm{CH}^{\cdot}(S)_{\mathbf{Q}}$, $\Psi_D(a * b)\,T^{2} =
\Psi_D(a)\,\Psi_D(b)\,T^{2}$ ; $T$ étant inversible,
\[
  \Psi_D(a * b) = \Psi_D(a)\,\Psi_D(b),
\]
sur toute base.\footnote{Une première rédaction de ce calcul, encadrée et
barrée d'une diagonale sur la page 9, part de $x * y = \mathrm{ch}(a * b)
T^{-1}$ et confond un instant $x$ et $\mathrm{ch}\,x$ ; la seconde, reprise
au-dessous, est celle donnée ici.} De même, en appliquant $\mathrm{ch}$ au
corollaire du théorème, $\varphi_L(xy) = (-1)^n\,\omega_A\,(\varphi_L(x) *
\varphi_L(y))$, et $(*)$ cette fois sur $B$ :
\[
  \Psi_D(ab) = (-1)^n\,\bigl(\Psi_D(a) * \Psi_D(b)\bigr)\;
  \mathrm{Todd}(t_A)\,\mathrm{Todd}(t_B)\,\mathrm{ch}(\omega_A).
\]
C'est ici que la lecture s'écarte de la page.\footnote{La page écrit un seul facteur $T$, avec « $T = \mathrm{Todd}(t_A)$,
$\omega_A = \hat\Lambda\,t_A^{\vee}$ » en accolade. Le calcul qu'elle fait
emploie $(*)$ pour le produit de Pontrjagin sur $B$, dont la classe de Todd
est $\mathrm{Todd}(t_B)$ ; il reste, après simplification par un $T$,
$\mathrm{Todd}(t_A)\,\mathrm{Todd}(t_B)\,\mathrm{ch}(\omega_A)$. Le terme de
degré $1$ de ce produit est $\frac{1}{2}c_1(t_A) + \frac{1}{2}c_1(t_B) -
c_1(t_A) = 0$, puisque $\tau_A \simeq \tau_B$ ; le produit vaut $1$ dès que
$S$ est local. Une note marginale, partiellement illisible, dit que $a, b \in
\mathrm{CH}^{\cdot}(A)_{\mathbf{Q}}$ et que « les formules \ldots dès que
\ldots ».} Sur un corps, les deux formules se réduisent à
\[
  \Psi_D(a * b) = \Psi_D(a)\,\Psi_D(b),
  \qquad
  \Psi_D(ab) = (-1)^n\,\Psi_D(a) * \Psi_D(b) :
\]
« ce sont les formules affirmées par Kleiman dans \emph{Dix exposés}, sur un
corps algébriquement clos et en travaillant avec une théorie cohomologique
plutôt qu'avec $\mathrm{CH}_{\mathbf{Q}}$ — à cela près que Kleiman n'a pas
tenu compte des facteurs $(-1)^n$ ! »\footnote{Le point d'exclamation est le
sien. Les formules avec le signe sont celles de Beauville (1983),
proposition 3 : $\mathcal{F}(x * y) = \mathcal{F}(x)\cdot\mathcal{F}(y)$ et
$\mathcal{F}(x \cdot y) = (-1)^g\,\mathcal{F}(x) * \mathcal{F}(y)$. Que
Kleiman ait omis le signe est l'affirmation de la page ; la lecture ne l'a
pas vérifiée.}

\section*{Remarques 4 et 5 : dualité et transport de structure (page 10)}

\subsection*{Remarque 4 : la transformée du dual}

Le morphisme $pr_2 \colon A \times B \to B$ est propre et lisse de dimension
relative $n$, de faisceau dualisant relatif $pr_1^{*}\omega_A[n]$ ; le
théorème de dualité donne donc, pour $F \in D_{\mathrm{parf}}(A)$ et
$F^{\vee}$ son dual,
\[
  \Phi_L(F)^{\vee} \;\simeq\; pr_{2*}\bigl(pr_1^{*}F^{\vee} \otimes L^{-1} \otimes pr_1^{*}\omega_A[n]\bigr)
  = \Phi_{L^{-1}}\bigl(F^{\vee} \otimes \omega_A[n]\bigr),
\]
et comme $L^{-1} \simeq (s_A \times \mathrm{id}_B)^{*}L$,
\[
  \Phi_L(F)^{\vee} \;\simeq\; \Phi_L\bigl(s_A^{*}F^{\vee} \otimes \omega_A[n]\bigr)
  \;\simeq\; s_B^{*}\,\Phi_L\bigl(F^{\vee} \otimes \omega_A[n]\bigr).
\]
Sur les groupes de Grothendieck :
$\varphi_L(x)^{\vee} = (-1)^n\,\omega_A\,\varphi_L(s_A^{*}x^{\vee})$.\footnote{La page écrit $\Phi_L(\check F \otimes \omega_A[n]) \simeq
\Phi_L(F)^{\vee}$ et $\varphi_L(x)^{\vee} = (-1)^n\,\omega_A\,
\varphi_L(\check x)$, sans la symétrie. Telle quelle la formule est fausse :
pour $F = e_s$, on a $\Phi_L(e_s)^{\vee} = L_s^{-1} = L_{-s}$, tandis que
$e_s^{\vee} \otimes \omega_A[n] \simeq e_s$ a pour transformée $L_s$. Le
dualisant de $L$ étant $L^{-1}$ et non $L$, la symétrie est nécessaire ; par
la remarque 5 elle peut être portée sur $A$ ou sur $B$. C'est la formule
(3.8) de Mukai, $\Phi(F^{\vee}) \simeq (-1)^{*}\Phi(F)^{\vee}[g]$. La lecture
corrige et signale.}

\subsection*{Remarque 5 : transport de structure}

Soit $u$ un automorphisme du schéma abélien $A$, et $u'$ l'automorphisme
\emph{contragrédient} de $B$, c'est-à-dire l'inverse du transposé de $u$,
caractérisé par $(u \times u')^{*}L \simeq L$. Par transport de structure,
\[
  \Phi_L(u_{*}F) \;\simeq\; u'_{*}\,\Phi_L(F),
\]
puisque $pr_1^{*}u_{*}F \otimes L \simeq (u \times \mathrm{id}_B)_{*}
\bigl(pr_1^{*}F \otimes (\mathrm{id}_A \times u'^{-1})^{*}L\bigr)$. En
particulier, $u = s_A$ ayant pour contragrédient $s_B$, la transformation
commute à la symétrie : $\Phi_L(s_A^{*}F) \simeq s_B^{*}\Phi_L(F)$. « Kif-kif
pour $\varphi_L$. »\footnote{La page écrit $\Phi_L(u.F) \simeq
u'.\Phi_L(F)$, le point désignant l'action ; elle ne dit pas lequel de $u^{*}$
ou $u_{*}$, qui sont inverses l'un de l'autre, et la lecture a choisi
$u_{*}$, ce qui fixe $u'$ comme l'inverse du transposé. Avec $u^{*}$ des deux
côtés la formule vaut également, pour le même $u'$.}

\end{document}
